\section{Related Work}\label{sec:related}

\subsection{Skepticism in long-term forecasting and component-level analysis}

A recurring finding in long-term time-series forecasting (LTSF) is that simple models rival elaborate architectures once evaluation is controlled. \citet{zeng2023dlinear} showed that one-layer linear maps match or beat contemporaneous Transformers; \citet{li2023rlinear} traced much of the remaining gap to a single component---reversible instance normalization---rather than to architectural depth; and \citet{toner2024linear} derived closed-form solutions for linear forecasters, showing that several normalization schemes amount to reparameterizations within the linear function class. Patch-based designs such as PatchTST \citep{nie2023patchtst} likewise attribute a large share of their gains to instance normalization. Benchmark-level audits reinforce this caution: TFB \citep{qiu2024tfb} documents how protocol details (e.g., dropping the last incomplete batch) distort comparisons, \citet{brigato2026champions} find no single champion across datasets, and \citet{position2025benchmark} argue that leaderboard-driven aggregation actively obscures \emph{when} and \emph{why} components help. Our study follows this component-level tradition, but replaces aggregate rankings with a conditional question: under what data-generating conditions does the normalization component itself help or hurt (\Cref{sec:theory})? Our variance-attribution analysis (\Cref{sec:empirical}) quantifies the stakes: in our grid, the normalization--dataset interaction explains far more error variance than any backbone-related factor.

\subsection{Normalization for nonstationary forecasting: the endogenous family}

RevIN \citep{kim2022revin} normalizes each input window by its own mean and variance and restores them at the output, targeting train--test distribution shift. Successors refine \emph{which endogenous statistics} to use and how to reinject them: Non-Stationary Transformers stationarize each window and restore the removed mean and variance through de-stationary attention \citep{liu2022nonstationary}, Dish-TS learns separate input and output distribution coefficients \citep{fan2023dishts}, SAN models slice-level statistics with a dedicated predictor \citep{liu2023san}, FAN removes dominant frequency components \citep{ye2024fan}, inner-instance normalization addresses point-level shift within a window \citep{inner2025instance}, IN-Flow learns the normalizing transformation itself via invertible flows \citep{inflow2024}, and SIN selects \emph{which} statistics to remove and restore by requiring them to be locally invariant yet globally variable \citep{han2024sin}. All of these estimate normalization statistics from the target series' own past---what we call the endogenous family. SIN is instructive for our purposes precisely because it poses a selection question and answers it inside that family: the candidate statistics are all functions of the lookback window, so no member of the search space can restore a level the window does not contain.

The heterogeneity of this family's effect is already visible at benchmark scale: \citet{zhang2024probts} report that instance normalization helps where trend dominates and fails on strongly seasonal, weakly trending series, tying the sign of the effect to data characteristics rather than to architecture.

Two recent papers re-examine the family critically, and one of them is the closest prior work to ours. \citet{role2026revin} ablate RevIN's components and report that some are redundant or even harmful, proposing a taxonomy of shift types to explain the heterogeneity. \citet{noise2025signal} deconstruct four contradictions in reversible normalization and take the nearest prior step toward the present contribution: they characterize each dataset by a portrait of endogenous statistics---among them a change-point risk $\mathrm{CPR}$ obtained by running a change-point detector over lookback windows---and pre-select a normalization for the whole dataset upfront by the rule ``if $\mathrm{CPR} < \tau$ use the robust variant, else use RevIN'', with $\tau$ a hand-set hyperparameter. Three differences separate their selector from the diagnostic we propose. First, both its inputs and its choice set are endogenous: it arbitrates between two instance-normalization variants using the target's own past, and so cannot express the decision we study, whereas $\lps$ scores the \emph{exogenous} predictability of the level and can therefore recommend restoring the level from covariates instead of from the lookback window. Second, their threshold is set by hand and left unanchored, whereas ours is fixed at the crossover implied by the closed-form model of \Cref{sec:theory} and committed before any run. Third---and this is their own headline finding---their selector suffers a systematic rather than random failure, and they close by arguing that a principled, diagnostics-driven strategy remains untapped. Our work supplies that version: a derived crossover (\Cref{sec:theory}), a pre-training score whose threshold the theory fixes, and an out-of-sample test of the resulting per-dataset predictions (\Cref{sec:lps,sec:empirical}). We emphasize that none of these distribution-shift normalization methods draws its normalization statistics from exogenous \emph{driver} series; the closest exception, APT \citep{li2026apt}, conditions only the affine rescaling step on deterministic calendar timestamps, leaving the instance mean and variance---and hence the level that is destroyed---estimated endogenously, and supplying no criterion for when instance normalization should be abandoned. The endogenous restriction is precisely the assumption our analysis relaxes. Two adjacent lines do use side information. \citet{chen2024solid} detect context-driven distribution shift and adapt a trained model at test time by fine-tuning on a contextually similar subset of the series' own past; this conditions on \emph{observed endogenous context} and adapts reactively at test time, whereas we condition on exogenous covariates and select the strategy before training. Closer to our normalizer, \citet{gamakumara2023conditional} normalize a series by a conditional mean and variance modeled as generalized additive functions of covariates---but as a statistical preprocessing step for imputation and correlation estimation, not as a normalization-selection question in multi-step forecasting, and without a boundary theory or a pre-training decision rule.

\subsection{Covariates in deep forecasters}

A parallel line incorporates exogenous information as \emph{model input}. iTransformer treats each variate, including covariates, as a token \citep{liu2024itransformer}; TimeXer dedicates cross-attention to exogenous series \citep{wang2024timexer}; CITRAS integrates covariates into a decoder-only Transformer \citep{citras2025}. These architectures decide how covariates enter the \emph{predictor}, whereas conditional normalization decides how covariates enter the \emph{normalization statistics}; the two mechanisms are orthogonal and composable. Indeed, our covariate-fair experiments (\Cref{sec:empirical}) grant every arm identical covariate inputs and still find that the normalization channel matters: instance normalization with covariates underperforms even a raw baseline with covariates, because the restore step overwrites level information the covariates carry.

\subsection{Feature-based forecasting-method selection}

Choosing a forecasting method from measurable properties of the series, rather than by running every candidate, is an established programme in the forecasting literature. FFORMS and FFORMA \citep{talagala2023metalearning, monteromanso2020fforma} compute a vector of time-series features---strength of trend and seasonality, autocorrelation structure, entropy, and so on---and learn a mapping from those features to a model choice or to combination weights, an approach whose value was demonstrated at scale in the M4 and M5 competitions \citep{makridakis2020m4, makridakis2022m5}. Our diagnostic belongs to this tradition and should be read within it: $\lps$ is a single feature, and the decision rule is a hand-specified threshold on it rather than a learned mapping. Three properties distinguish it from the meta-learning approach. The feature is not endogenous---every feature in the standard FFORMS set is computed from the target series alone, whereas $\lps$ is by construction a property of the target \emph{jointly with} its covariates. The target of the decision is a \emph{component} of a pipeline rather than a whole method. And the threshold is derived from a risk comparison (\Cref{sec:theory}) rather than fitted, which is what makes an out-of-sample test of its predictions meaningful with only eight series; a learned mapping would need far more. The cost of that choice is equally clear: a one-feature, hand-thresholded rule cannot express interactions that a meta-learner would capture, and \Cref{sec:empirical} reports a case where it does not order what it might be expected to.

\subsection{Energy forecasting practice}

Practitioners have long normalized by \emph{external} or physically grounded quantities: wind power is divided by installed capacity, solar output by a clear-sky index, and load is regressed on temperature-derived degree days---practices standardized in the GEFCom2014 competitions \citep{hong2016gefcom}---and adaptive normalization schemes are studied for wind data specifically \citep{patil2022wind}. These are, in our terms, implicit forms of conditional normalization: the normalizer is driven by covariates rather than by the target's recent past. This conditioning was not incidental: it encodes the domain knowledge that the level of a renewable or load series is largely set by exogenous drivers---exactly the information the endogenous instance-normalization default discards. What practice lacked was not the instinct but a \emph{criterion}: when is such conditioning necessary rather than optional---i.e., when does the endogenous default actively destroy forecastable signal? Our contribution supplies it: the crossover condition of \Cref{sec:theory} and the LPS rule of \Cref{sec:lps} tell a practitioner, before training, whether capacity- or weather-driven normalization is optional polish or a requirement.

\subsection{Classical dynamic regression}

Conditioning a forecast on exogenous drivers long predates deep forecasting. Box--Jenkins transfer-function models and regression with ARIMA errors \citep{box1970time, hyndman2021fpp} regress the target on exogenous covariates and model the residual autocorrelation with an ARIMA process: covariates set the conditional level and a stochastic term captures the rest. In energy this lineage is the operational standard---probabilistic wind and load forecasts are built on covariate-driven models with deliberate treatment of the exogenous--stochastic split \citep{pinson2013wind, hong2016gefcom}. Instance normalization sits orthogonally to, and in the exogenous regime in conflict with, this tradition: by estimating level statistics from the target's own past window it removes precisely the level a dynamic regression would have explained from covariates. Accordingly, a fair evaluation must compare against a properly specified regression-with-ARIMA-errors model rather than a strawman autoregressive one---a claim that classical covariate conditioning was right cannot rest on a weak classical implementation.

\paragraph{Gap} Summarizing across these strands: component-level analyses identify instance normalization as pivotal but rank it unconditionally; the endogenous normalization family and its recent critics document failures without a predictive theory of them; covariate-aware architectures use exogenous information as input but not as normalization statistics; and both energy practice and classical dynamic regression condition on covariates---the latter formally, the former implicitly---but neither supplies a criterion for when the endogenous deep-learning default destroys the signal that conditioning would preserve. Individual ingredients have precedents---covariate-conditional normalization as a statistical preprocessing step \citep{gamakumara2023conditional}, context-driven test-time adaptation \citep{chen2024solid}, empirical critiques of reversible normalization \citep{role2026revin}, an endogenous, hand-thresholded pre-training normalization selector \citep{noise2025signal}, statistic selection inside the endogenous family \citep{han2024sin}, black-box meta-learning over design choices \citep{liang2025tsgym}, and feature-based forecasting-method selection \citep{talagala2023metalearning, monteromanso2020fforma}. What none of them supplies, and what this paper contributes, is the combination of (i) a risk comparison that says when instance normalization hurts and by how much, (ii) a contrast between \emph{exogenous} and endogenous normalization statistics---so that the choice set includes not estimating the level from the window at all---and (iii) a decision rule on a single pre-training quantity whose threshold is derived from (i) rather than learned or hand-set, with its per-dataset predictions committed before the experiments that test them.
